\mwnewpart
\renewcommand\runningtitle{Section 6. Install}
\chapter{Install}
\sectioncontents

\section{What Install Does}\ixs{Install}

\begin{plainwords}
PC DOS 2.10 had no installation program: the owner of an XT partitioned
the disk, formatted it and copied the files by hand, following the
manual. RetrOS/LT adds Install, as every RetrOS edition has. It asks
before each change, works out what has already been done from the state
of the disk, and so can simply be run again after the restart a new
partition needs.
\end{plainwords}

\file{install/INSTALL.ASM} is the RetrOS/x86 Install adapted to DOS 2.
The steps (Section~2 of the \textit{Guide to Operations} shows the
screens):

\begin{enumerate}
\item The license, \file{LICENSE.TXT}, must be accepted.
\item The first fixed disk's master boot record is read; with no type~1
  partition, Install writes one over the whole disk, or runs
  \typed{FDISK}, and restarts the computer.
\item The directory is asked for, \file{C:\bs{}DOS} unless changed.
\item An unformatted partition gets \typed{FORMAT C:/S}; a formatted one
  \typed{FORMAT} or \typed{SYS C:}, as the user chooses. Either way
  Install then checks that \file{C:\bs{}RETBIO.COM} is there.
\item The DOS diskette's files are copied, then the Supplemental Programs
  diskette's, each checked by its \file{RETROS.DSK}.
\item \file{CONFIG.SYS} and \file{AUTOEXEC.BAT} are written, the old ones
  kept as \file{.OLD}.
\end{enumerate}

\section{What Changed for DOS 2}

\subsection{Running a program} DOS 2 programs take their drives from the
two file control blocks in the program segment prefix, which
\file{COMMAND.COM} fills from the command line; the RetrOS/x86 Install,
for DOS 4, left them empty, and \typed{FORMAT} took drive A. Install now
fills them itself before EXEC, as \file{COMMAND.COM} does:

\srclist{install_fcb}{The file control blocks, from the command tail.}

\engnote{The first install test caught this: \typed{FORMAT} asked for a
diskette in drive A. Nothing was harmed, since \typed{FORMAT} waits for a
key before formatting a diskette; but it is why Install now checks for
\file{RETBIO.COM} after \typed{FORMAT} and \typed{SYS}, whose exit codes
in DOS 2 say nothing.}

\subsection{The partition} RetrOS/x86's Install ran \typed{FDISK} with
switches; IBM's \typed{FDISK} for DOS 2, and RetrOS/LT's, has none. Install
writes the whole-disk partition itself, \routine{whole\_disk}, exactly as
\typed{FDISK} makes it: active, type~1, from head~1 of cylinder~0, at most
65,535 sectors, with RetrOS/LT's master boot record code.

\subsection{Errors} DOS 2's critical error handler has no ``Fail'' answer,
so Install leaves DOS's own handler in place.

\subsection{Disk transfers} Install reads the master boot record and the
partition's startup record through the ROM, into its 16 KB buffer, and
writes the master boot record from it. A sector there could cross a 64 KB
boundary, so the transfer uses whichever half of the buffer's first 1,024
bytes does not:

\srclist{install_dma}{A buffer the DMA controller can use.}

\subsection{The startup files} \file{CONFIG.SYS} has no \code{SHELL=}
(\file{COMMAND.COM} is in the root, where DOS 2 looks), \code{FILES=20},
\code{BUFFERS} by memory, and \file{ANSI.SYS}. \file{AUTOEXEC.BAT} asks
the date and time, since neither the PC nor the XT keeps them, sets
\code{PROMPT \$P\$G} and the path.

\section{The Diskettes}\ixs{diskettes}\ixs{mkinstall.py}

\file{scripts/mkinstall.py} makes the two 180 KB images: a blank
single-sided diskette with RetrOS/LT's startup record and a media byte of
FCh, then the files in IBM's order. \file{RETBIO.COM} and
\file{RETDOS.COM} go first, so that they are the first two directory
entries and \file{RETBIO.COM} is in one piece at the start of the data
area. Every file is given the release date and time, as IBM's diskettes
had one date for all; the hidden \file{RETROS.DSK} names the diskette for
Install. The DOS diskette's \file{AUTOEXEC.BAT} is \typed{DATE},
\typed{TIME}, \typed{INSTALL}.
